% stack_tecnico.tex — Capítulo 2: Stack Tecnológico
\chapter{Stack Tecnológico}

\section{Runtime y Framework Principal}

\subsection{Node.js}

Node.js es un entorno de ejecución de JavaScript del lado del servidor construido sobre el motor V8 de Google Chrome. Para el backend de FINARA se utilizó la versión \textbf{18.x LTS} (Long-Term Support), garantizando estabilidad y soporte de seguridad a largo plazo \cite{nodejs_docs}.

\begin{notaTecnica}
  Node.js implementa un modelo de I/O no bloqueante basado en el \textit{event loop}, lo cual lo hace especialmente eficiente para APIs REST donde las operaciones dominantes son lecturas y escrituras a base de datos, no cómputo intensivo \cite{nodejs_event_loop}.
\end{notaTecnica}

\subsection{Express.js}

Express.js 4.18.x es el framework web minimalista utilizado sobre Node.js \cite{express_docs}. Proporciona el enrutamiento HTTP, el sistema de middleware y el manejo de errores del backend FINARA.

Los middleware de Express utilizados en el proyecto se muestran en la Tabla \ref{tab:middleware}.

\begin{longtable}{p{4cm} p{9cm}}
  \caption{Middleware Express configurados en app.js}
  \label{tab:middleware} \\
  \toprule
  \textbf{Middleware} & \textbf{Función} \\
  \midrule
  \endfirsthead
  \toprule
  \textbf{Middleware} & \textbf{Función} \\
  \midrule
  \endhead
  \texttt{cors()} & Habilita solicitudes Cross-Origin desde la app Flutter \\
  \texttt{express.json()} & Parsea body de requests JSON (límite 10MB) \\
  \texttt{express.urlencoded()} & Parsea body de formularios codificados \\
  \texttt{express.static()} & Sirve imágenes de tickets desde \texttt{/uploads} \\
  \texttt{logger} & Log coloreado de cada request con status code y tiempo \\
  \texttt{sanitizer} & Escapado de caracteres HTML en todos los inputs \\
  \texttt{apiLimiter} & Rate limiting: 200 req/15min para la API general \\
  \texttt{authLimiter} & Rate limiting: 10 req/15min en endpoints de auth \\
  \texttt{ocrLimiter} & Rate limiting: 20 req/15min en escaneo de tickets \\
  \bottomrule
\end{longtable}

\section{Base de Datos}

\subsection{MySQL 8.0}

MySQL 8.0 con motor de almacenamiento InnoDB es el sistema de gestión de base de datos relacional utilizado \cite{mysql_docs}. InnoDB proporciona soporte de transacciones ACID, claves foráneas y bloqueo a nivel de fila, características esenciales para mantener la integridad de los datos financieros del usuario \cite{mysql_innodb}.

El driver de conexión utilizado es \textbf{mysql2} con el módulo \texttt{promise}, que expone una API basada en \texttt{async/await} y utiliza \textit{prepared statements} de forma nativa, previniendo inyecciones SQL \cite{mysql2_npm}.

\begin{lstlisting}[style=javascript, caption={Configuración del pool de conexiones MySQL}, label={lst:db-pool}]
// src/config/database.js
const mysql = require('mysql2/promise');

const pool = mysql.createPool({
  host:               process.env.DB_HOST,
  port:               parseInt(process.env.DB_PORT) || 3306,
  user:               process.env.DB_USER,
  password:           process.env.DB_PASSWORD,
  database:           process.env.DB_NAME,
  waitForConnections: true,
  connectionLimit:    10,     // Maximo 10 conexiones simultaneas
  timezone:           'local',
  enableKeepAlive:    true,   // Reconexion automatica
  keepAliveInitialDelay: 30000,
});
\end{lstlisting}

\section{Dependencias de Producción}

La Tabla \ref{tab:deps-prod} describe todos los paquetes npm declarados en \texttt{dependencies} del \texttt{package.json}.

\begin{longtable}{>{\ttfamily}p{3.8cm} >{\centering}p{2cm} p{7.5cm}}
  \caption{Dependencias de producción del backend FINARA}
  \label{tab:deps-prod} \\
  \toprule
  \textbf{Paquete} & \textbf{Versión} & \textbf{Justificación Técnica} \\
  \midrule
  \endfirsthead
  \toprule
  \textbf{Paquete} & \textbf{Versión} & \textbf{Justificación Técnica} \\
  \midrule
  \endhead
  express            & 4.18.3  & Framework web minimalista y maduro; ecosistema de middleware extenso \cite{express_docs} \\
  \addlinespace
  mysql2             & 3.9.2   & Driver MySQL con prepared statements nativos y API Promise; más rápido que el paquete \texttt{mysql} original \cite{mysql2_npm} \\
  \addlinespace
  bcrypt             & 5.1.1   & Algoritmo de hash con sal adaptativa para contraseñas; resistente a ataques de fuerza bruta \cite{bcrypt_npm} \\
  \addlinespace
  jsonwebtoken       & 9.0.2   & Implementación del estándar RFC 7519 (JWT) para autenticación sin estado \cite{jwt_rfc} \\
  \addlinespace
  cors               & 2.8.5   & Middleware CORS con configuración de orígenes permitidos por variable de entorno \\
  \addlinespace
  dotenv             & 16.4.5  & Carga de variables de entorno desde \texttt{.env}; evita credenciales en código fuente \cite{dotenv_npm} \\
  \addlinespace
  multer             & 1.4.5-lts.1 & Middleware multipart/form-data para recepción de imágenes de tickets; límite configurable \\
  \addlinespace
  express-validator  & 7.0.1   & Librería de validación y sanitización de inputs basada en validator.js \cite{express_validator} \\
  \addlinespace
  uuid               & 9.0.1   & Generación de identificadores únicos universales v4 para nombrado de archivos subidos \\
  \bottomrule
\end{longtable}

\section{Dependencias de Desarrollo}

\begin{longtable}{>{\ttfamily}p{3.8cm} >{\centering}p{2cm} p{7.5cm}}
  \caption{Dependencias de desarrollo del backend FINARA}
  \label{tab:deps-dev} \\
  \toprule
  \textbf{Paquete} & \textbf{Versión} & \textbf{Propósito} \\
  \midrule
  nodemon   & 3.1.0  & Reinicio automático del servidor al detectar cambios en archivos durante desarrollo \\
  \addlinespace
  jest      & 29.7.0 & Framework de testing con soporte de mocks, cobertura de código y modo watch \cite{jest_docs} \\
  \addlinespace
  supertest & 6.3.4  & Pruebas de integración HTTP sobre Express sin levantar un servidor real en red \cite{supertest_npm} \\
  \bottomrule
\end{longtable}

\section{package.json Completo}

\begin{lstlisting}[style=json, caption={package.json del backend FINARA v4.0.0}, label={lst:package-json}]
{
  "name": "finara-backend",
  "version": "4.0.0",
  "description": "API REST - FINARA App de administracion financiera",
  "main": "server.js",
  "scripts": {
    "start":       "node server.js",
    "dev":         "nodemon server.js",
    "test":        "jest --coverage --detectOpenHandles",
    "test:watch":  "jest --watch",
    "test:single": "jest --testPathPattern"
  },
  "jest": {
    "testEnvironment": "node",
    "testMatch": ["**/tests/**/*.test.js"],
    "coverageThreshold": {
      "global": { "functions": 60, "lines": 60 }
    }
  },
  "dependencies": {
    "bcrypt": "^5.1.1",
    "cors": "^2.8.5",
    "dotenv": "^16.4.5",
    "express": "^4.18.3",
    "express-validator": "^7.0.1",
    "jsonwebtoken": "^9.0.2",
    "multer": "^1.4.5-lts.1",
    "mysql2": "^3.9.2",
    "uuid": "^9.0.1"
  },
  "devDependencies": {
    "jest": "^29.7.0",
    "nodemon": "^3.1.0",
    "supertest": "^6.3.4"
  }
}
\end{lstlisting}

\section{Stack OCR (Python)}

El módulo de reconocimiento óptico de caracteres es un componente externo escrito en Python 3.x, invocado desde Node.js mediante \texttt{child\_process.spawn}. Las dependencias Python se detallan en la Tabla \ref{tab:deps-ocr}.

\begin{longtable}{>{\ttfamily}p{3.5cm} p{9.5cm}}
  \caption{Dependencias Python del módulo OCR}
  \label{tab:deps-ocr} \\
  \toprule
  \textbf{Librería} & \textbf{Función} \\
  \midrule
  opencv-python & Preprocesamiento de imágenes: upscaling adaptativo, filtro bilateral, corrección de iluminación, binarización y sharpening \cite{opencv_docs} \\
  \addlinespace
  pytesseract   & Wrapper Python para el motor Tesseract OCR de Google; extrae texto con coordenadas y confianza por palabra \cite{tesseract_docs} \\
  \addlinespace
  numpy         & Operaciones matriciales sobre imágenes: kernels de convolución, manipulación de píxeles \cite{numpy_docs} \\
  \addlinespace
  difflib       & Módulo estándar de Python; usado para fuzzy matching de tiendas sin dependencias externas \\
  \bottomrule
\end{longtable}

\section{Herramientas de Desarrollo}

\begin{itemize}[leftmargin=1.5cm]
  \item \textbf{Visual Studio Code}: editor principal con extensiones ESLint, Prettier y REST Client.
  \item \textbf{Postman}: colecciones de prueba de los 46 endpoints con variables de entorno y scripts de test automáticos.
  \item \textbf{MySQL Workbench}: diseño del esquema ER y ejecución del script \texttt{schema.sql}.
  \item \textbf{Git}: control de versiones; repositorio privado en GitHub.
  \item \textbf{Overleaf}: compilación y edición de esta documentación LaTeX.
\end{itemize}
